% ======================================================================
\section{Method}
\label{sec:method}
% ======================================================================

We use two model names to avoid temporal-stage ambiguity. The
\emph{geometry-only policy} is the port-Hamiltonian navigator trained
without context risk;
the \emph{context-enriched field} is the same update after adding
one context-energy term and, when used, the local affordance gate.
Thus the enrichment is an architectural and training extension, not
the second temporal stage of an online rollout. The added term introduces
a learned force channel whose gradient reshapes the closed-loop decision
field only when the local scene admits a lower-risk feasible maneuver.  At each step the agent observes goal
displacement $\vec c_g-\vec c_t$, obstacle features within radius $\hat d$,
and a $P\times P$ local context patch $P_t$ containing a smoothed soft-risk
field $\tilde r(\cdot)\in[0,1]$ and a signed-distance field $\phi(\cdot)$ to
hard hazards.  These fields are bilinearly resampled at the current position
during integration.  All baselines and ablations use the same local-information
contract.

% ----------------------------------------------------------------------
\subsection{Shared Hamiltonian update}
\label{sec:method:skeleton}
% ----------------------------------------------------------------------

The geometry-only policy and the context-enriched field share the same phase-space update; the only dynamical
addition is the boldfaced context-force term in the momentum equation:
\begin{equation}
\renewcommand{\arraystretch}{1.5}
\begin{array}{ll}
\textbf{Geometry-only policy} & \textbf{Ctx-enriched} \\[2pt]
\Mom_{t+1} = \Mom_t
  - \tau\nabla_\Pos\Ham_{\rm geom}
  - \tau R_\theta(\Mom_t)
&
\Mom_{t+1} = \Mom_t
  - \tau\nabla_\Pos\Ham_{\rm geom}
  \;\boldsymbol{-\tau\nabla_\Pos\Ham_{\rm ctx}}
  - \tau R_\theta(\Mom_t) \\[4pt]
\multicolumn{2}{l}{\Pos_{t+1} = \Pos_t + \tau\Mass^{-1}\Mom_{t+1}
  \quad\text{(identical in both fields).}}
\end{array}
\label{eq:method:skeleton}
\end{equation}
New sensory information therefore enters through the cotangent variable
$\Mom$, not through a separate critic or global replanner.  The decision field
is the policy: changing the stored energy changes the vector field that
generates the next motion.

% ----------------------------------------------------------------------
\subsection{Context energy and route-aware force channel}
\label{sec:method:energy}
% ----------------------------------------------------------------------

The context term is additive, so the context-enriched field preserves the geometry-only policy
when its context coefficients vanish:
\begin{equation}
\Ham_\theta(\Pos,\Mom)
=
\underbrace{
  \tfrac{1}{2}\Mom^\top\Mass^{-1}\Mom
  + \beta D_{\mathcal{G}}(\Pos,\vec c_g)
  + \textstyle\sum_j\alpha_j b_{\rm IPC}(d_j)
}_{\Ham_{\rm kin}+\Ham_{\rm geom}\;\text{(geometry-only)}}
+
\underbrace{
  \lambda_s\tilde r(\vec c)
  + \lambda_h b_{\rm sp}(\phi(\vec c))
}_{\Ham_{\rm ctx}\;\text{(Ctx-enriched only)}}.
\label{eq:method:energy}
\end{equation}
Here $b_{\rm IPC}$~\citep{li2020ipc} is the incremental-potential contact
barrier and $b_{\rm sp}(\phi)=k^{-1}\log(1+e^{k(\hat d_\phi-\phi)})$ is a
softplus-relaxed inverse SDF.  Differentiating $\Ham_{\rm ctx}$ yields two
force channels:
\begin{equation}
F_{\rm ctx}(\vec c)
=
\underbrace{-\lambda_s\nabla\tilde r(\vec c)}_{F_{\rm soft}:\;\text{lateral risk deflection}}
+
\underbrace{-\lambda_h b_{\rm sp}'(\phi(\vec c))\nabla\phi(\vec c)}_{F_{\rm hard}:\;\text{hazard repulsion}},
\qquad
b_{\rm sp}'(\phi)=-\sigma(k(\hat d_\phi-\phi)).
\label{eq:method:force}
\end{equation}
$F_{\rm soft}$ expresses preferences among feasible maneuvers, while
$F_{\rm hard}$ is a differentiable penalty against boundary contact. It is not
a control-barrier certificate; throughout the paper, violation metrics should
be read as empirical tail-risk reduction under finite rollouts. Risk is
therefore not a post-hoc cost, but a force that reshapes the local closed-loop
dynamics.

\begin{figure}[t]
  \centering
  \includegraphics[width=0.95\linewidth]{NeurIPS_2026/figures/ch4_factored_hamiltonian.png}
  \caption{\textbf{Factored stored energy and induced force channels.}
  $\Ham_\theta$ separates kinetic, geometric, dissipative, and context
  terms. The context term creates a soft-risk deflection channel and a
  hard-hazard repulsion channel. The route-aware gate lets the soft channel
  shift the field only when a feasible lower-risk maneuver exists; otherwise
  the rollout stays near the geometry-only policy.}
  \label{fig:factored}
\end{figure}

The scalar soft-risk force can overreact when a low-risk-looking region is
blocked by a fence, berm, vehicle, or other non-traversable boundary.  We
therefore gate only the soft-risk channel using a local affordance test:
\begin{equation}
F_{\rm ctx}^{\rm ra}(\vec c;P_t)
=
m_{\rm feas}(\vec c,z_t)\bigl(-\lambda_s\nabla\tilde r(\vec c)\bigr)
-
\lambda_h b_{\rm sp}'(\phi(\vec c))\nabla\phi(\vec c),
\label{eq:method:route_aware_force}
\end{equation}
where $z_t$ denotes the finite set of $K$ short-horizon primitives sampled
inside the current BEV patch.  Let $k^\star$ be the feasible primitive with
lowest soft-risk exposure, and let $k_{\rm geo}$ be the primitive closest to
the geometry-only policy direction.  The gate is
\begin{equation}
m_{\rm feas}(\vec c,z_t)
=
\sigma\!\left(\kappa_R[(R_{k_{\rm geo}}-R_{k^\star})-\rho_R]\right)
\sigma\!\left(\kappa_C(C_{k^\star}-\delta_\phi)\right)
I_{k^\star}.
\label{eq:method:route_gate}
\end{equation}
Here $R_k$ is the primitive's soft-risk integral, $C_k$ its minimum SDF
clearance, and $I_k$ its local traversability indicator.  Thus
$F_{\rm soft}$ activates only when the current patch contains a feasible
candidate that improves risk by margin $\rho_R$; otherwise it is suppressed
without weakening $F_{\rm hard}$.  The gate is not a global replanner: it only
tests local affordance inside the same sensed patch available to all
local-sensing methods. Under semantic or risk-map corruption, the soft channel
therefore has a conservative failure mode: a noisy low-risk patch is not enough
to activate lateral deflection unless a short-horizon primitive also clears the
SDF and improves risk by margin $\rho_R$. This is not a perception guarantee,
but it tends to turn uncertain or inconsistent risk evidence into suppression
rather than an unconstrained evasive force.  Figure~\ref{fig:route_gate_main}
summarizes the gate used in all experiments.  Coefficients $(\lambda_s,\lambda_h)$ are predicted
by sigmoid-bounded CNN heads from $P_t$ fused with the goal token; geometric
heads $(\alpha,\beta,\gamma)$ are frozen from the geometry-only policy, so these coefficients
are the only route by which context alters the field
(Appendix~\ref{app:coeff_pred}).

\begin{figure}[t]
\centering
\fbox{\begin{minipage}{0.94\linewidth}
\small
\textbf{Route-aware soft-channel gate.}
Given the current BEV patch, sample short-horizon primitives $z_t$ and compute
each primitive's soft-risk integral $R_k$, minimum SDF clearance $C_k$, and
traversability flag $I_k$. Choose the lowest-risk feasible primitive $k^\star$
and compare it with the primitive $k_{\rm geo}$ closest to the geometry-only
policy direction. The soft channel opens only if
\[
R_{k_{\rm geo}}-R_{k^\star}>\rho_R,\qquad
C_{k^\star}>\delta_\phi,\qquad I_{k^\star}=1.
\]
If any test fails, $m_{\rm feas}\approx0$: $F_{\rm soft}$ is suppressed while
$F_{\rm hard}$ remains active. The test is local, differentiable through the
sigmoids in Eq.~\eqref{eq:method:route_gate}, and does not call a global planner.
\end{minipage}}
\caption{\textbf{Gate specification in the main method.}
The gate converts a risk-map cue into force activation only when the current
local patch contains a cleared, traversable primitive that improves soft risk
by margin $\rho_R$.}
\label{fig:route_gate_main}
\end{figure}

% ----------------------------------------------------------------------
\subsection{Tail-risk objective}
\label{sec:method:cvar}
% ----------------------------------------------------------------------

Each rollout accumulates
\begin{equation}
J(\theta) =
w_g\|\Pos_T-\Pos_g\|^2
+ w_\ell\textstyle\sum_t\|\Pos_{t+1}-\Pos_t\|
+ w_r\textstyle\sum_t\tilde r(\Pos_t)\|\Pos_{t+1}-\Pos_t\|
+ w_h\textstyle\sum_t\mathbf{1}[\phi(\Pos_t)<\epsilon].
\label{eq:method:J}
\end{equation}
Because the relevant failures are rare, expected-cost training can be dominated
by typical rollouts.  We therefore optimize the empirical
Rockafellar--Uryasev CVaR objective~\citep{rockafellar2000optimization}:
\begin{equation}
\widehat{\cvar{\alpha}}(\theta)
=
\hat\eta
+\frac{1}{(1-\alpha)B}
\sum_{i=1}^{B}\bigl(J^{(i)}(\theta)-\hat\eta\bigr)_+,
\qquad
\hat\eta=\widehat Q_\alpha\!\bigl(\{J^{(i)}\}_{i=1}^B\bigr)\;\text{detached}.
\label{eq:method:cvar_emp}
\end{equation}
With $\alpha=0.95$ and $B=64$, gradient flows through the worst rollouts,
giving a consistent estimator~\citep{hong2014monte}.  The resulting gradient
acts through the sensitivity
$\partial \Pos_{t+1}/\partial(\lambda_s,\lambda_h)$ exactly on rollouts where
the context force can change the outcome.

% ----------------------------------------------------------------------
\subsection{Training hierarchy}
\label{sec:method:loops}
% ----------------------------------------------------------------------

The context-enriched field uses three computational roles.  First, the geometry-only checkpoint is
loaded with $\lambda_s,\lambda_h\equiv0$, so Eq.~\eqref{eq:method:skeleton}
initially reproduces the geometry-only policy.  Second, the episode loop
backpropagates the CVaR loss in Eq.~\eqref{eq:method:cvar_emp} through the
differentiable rollout to update the context encoder and coefficient heads.
Third, a lightweight segment-correction loop updates active geometry/context
coefficients from realized local residuals, while a curriculum loop advances
to harder patches only after validation violation rates fall below a threshold.
The full loop diagram, initialization protocol, and pseudocode are given in
Appendix~\ref{app:ctx_loops_hierarchy}; the main text keeps the algorithmic
contract to the force law, loss, and update hierarchy above.
For the logged $500$-epoch highway runs, the geometry-only policy trained for
$1.49$ hours and the context-enriched policy for $3.27$ hours on one NVIDIA
A100 GPU, for $4.8$ total GPU-hours.

% ----------------------------------------------------------------------
\subsection{Selectivity consequences of context-energy enrichment}
\label{sec:method:enrichment}
% ----------------------------------------------------------------------

The additions above are coupled: adding $\Ham_{\rm ctx}$ creates the force
channel $F_{\rm ctx}=-\nabla_\Pos\Ham_{\rm ctx}$, the coefficient subspace
$(\lambda_s,\lambda_h)$, and the rollout-sensitivity pathway through which
CVaR training acts.  The same structure yields the selectivity consequences
measured in Sec.~\ref{sec:experiments}.

\begin{proposition}[Selectivity consequences of context-energy enrichment]
\label{prop:enrichment}
Let the geometry-only policy and the context-enriched field share the same geometric Hamiltonian, dissipative
port $R_\theta$, mass matrix $\Mass$, and semi-implicit integration rule.  Let
$\Ham_{\rm enr}=\Ham_{\rm geo}+\Ham_{\rm ctx}$ with
$F_{\rm ctx}=-\nabla_\Pos\Ham_{\rm ctx}$.  Under the boundedness and local
Lipschitz assumptions in Appendix~\ref{app:theory}, the following hold over
any finite horizon:
\begin{enumerate}
  \item[\textbf{C1.}] \emph{Geometry-only preservation.} If
  $\sup_t\|F_{\rm ctx}(q_t)\|\le\varepsilon$, then the context-enriched rollout remains
  within $O(\varepsilon)$ of the geometry-only rollout.
  \item[\textbf{C2.}] \emph{No hallucinated escape.} If
  $\sup_t\|P_\perp F_{\rm ctx}(q_t)\|\le\varepsilon_\perp$, then lateral
  deviation from the geometry-only rollout is bounded by $O(\varepsilon_\perp)$.
  \item[\textbf{C3.}] \emph{Selective risk deflection.} If a feasible lateral
  direction has projected risk-gradient margin $\Delta$ and projected
  hard-barrier forces cancel the soft-risk force by at most $\chi$, then
  $\|P_\perp F_{\rm ctx}\|\ge\lambda_s\Delta-\chi$ and one semi-implicit step
  decreases local soft risk relative to the geometry-only policy by
  $O(\tau^2\lambda_s\Delta^2)$ whenever $\lambda_s\Delta>\chi$.
\end{enumerate}
\end{proposition}

Proofs: C1 $\Rightarrow$ Theorem~\ref{thm:scaffold_preservation};
C2 $\Rightarrow$ Corollary~\ref{cor:no_hallucination};
C3 $\Rightarrow$ Theorem~\ref{thm:risk_deflection}.  Selectivity is not
trained as a separate label: C1--C3 follow from the gradient structure of
$\Ham_{\rm ctx}$ and the affordance gate, not from post-hoc tuning. Appendix
Table~\ref{tab:app_theory_checks} makes C1--C3 checkable by mapping each
theoretical consequence to measured rollout statistics.

\begin{table}[t]
\centering
\caption{\textbf{Main-text check of the theory predictions.}
We test C1--C3 as finite-rollout predictions, not formal safety certificates.
``Below-$\delta$'' means lateral deviation stays below the reaction threshold
used for false pre-activation.}
\label{tab:main_theory_checks}
\small
\setlength{\tabcolsep}{4pt}
\renewcommand{\arraystretch}{1.06}
\begin{tabular}{lccc}
\toprule
Prediction & Rollout condition & Below-$\delta \uparrow$ & Failure metric$\downarrow$ \\
\midrule
C1 & geometry-only delayed escape
& $1.000$ & false pre-act $0.000$ \\
C2 & route-aware context, static blocked R2
& $0.886$ & FAR $0.114$ \\
C3 & route-aware CVaR, pre-escape delayed
& $0.820$ & false pre-act $0.180$ \\
\bottomrule
\end{tabular}
\end{table}
